\section{Introduction}
\label{sec:introduction}

Large language models (LLMs) pre-trained on vast corpora have become the dominant
paradigm for natural language processing \citep{Vaswani2017, Devlin2019}. A single
pre-trained transformer backbone can be adapted to dozens of downstream tasks through
fine-tuning, and parameter-efficient fine-tuning (PEFT) methods have reduced the cost
of this adaptation to a fraction of the original model size \citep{Houlsby2019, Hu2022}.
The result is an ecosystem in which a single frozen backbone serves as a universal
substrate, and small task-specific adapter modules provide the task-particular
knowledge required for each application.

This ecosystem creates an unprecedented opportunity for collaborative knowledge sharing.
If adapters are small, self-contained, and additive to a shared frozen backbone, they
are in principle \textit{transmissible}: a model trained by one party on a specialist
task can be shared with another party who needs that capability, without revealing
the underlying training data, and without requiring the receiving party to retrain the
backbone from scratch. In a peer-to-peer (P2P) network of nodes sharing the same
frozen backbone, adapters are, in principle, transmissible units of knowledge exchange. Each peer
may contribute specialist adapters it has trained, and benefit from adapters
contributed by others.

\textbf{The research question} addressed by this review is: \textit{is a
peer-to-peer research direction for adapter-based multi-task inference
theoretically grounded and supported by the current state of the literature?}
More specifically: given a network of autonomous nodes each holding a frozen
transformer backbone and a local collection of task-specific adapters, does the
existing literature provide the building blocks required for any node to
discover, acquire, compose, and serve adapters from across the network for
arbitrary multi-task NLP queries, without any central coordinator, and if
not, what are the open gaps?
Section~\ref{sec:methodology} operationalises this feasibility question into a literature-searchable form: \textit{what methods have been proposed 
for decentralised, adapter-based multi-task inference over a shared frozen 
transformer backbone, and what are the open research gaps?}

This problem sits at the intersection of three active research fields:
\begin{enumerate}
  \item \textbf{PEFT and adapter architectures} (Section~\ref{sec:peft},
    Section~\ref{sec:adapter_composition}): What are the structural properties of
    adapter modules that make them suitable as knowledge exchange units?
  \item \textbf{Efficient inference and routing} (Section~\ref{sec:inference_systems},
    Section~\ref{sec:moe_routing}): How can a serving system select and compose the
    right adapter for each request at minimal latency?
  \item \textbf{Distributed systems and federated learning} (Section~\ref{sec:p2p_federated}):
    What algorithms and protocols enable model parameter exchange without a central
    server, and what are the convergence guarantees under heterogeneous data
    distributions?
\end{enumerate}

The multi-tenant serving literature (S-LoRA \citep{Sheng2024}, Punica \citep{Chen2023Punica}, 
CaraServe \citep{Li2024CaraServe}) assumes a single administrative domain. 
The federated PEFT literature assumes a central server for adapter aggregation. 
The P2P inference literature (Petals \citep{Borzunov2023}) distributes backbone 
layers across nodes but does not address per-task adapter modules. The closest precedent is 
\citet{Sajina2024}, which demonstrates that task-specific modules can be trained and served 
across a P2P network of encoder-only transformers — but the modules are tied to a fixed peer 
population, are not framed as PEFT/LoRA-style adapters, and are neither discoverable nor composable 
by peers outside the original training group. Each of these bodies of work addresses part
of the problem, but their intersection has not yet been addressed as an integrated whole.

\subsection{Review Structure}
\label{sec:introduction:structure}

Section~\ref{sec:background} provides the technical background on transformer
architectures, PEFT fundamentals, and P2P network primitives. Section~\ref{sec:methodology}
describes the systematic literature review methodology. Sections~\ref{sec:peft}--\ref{sec:p2p_federated}
review the five bodies of related work. Section~\ref{sec:synthesis_gap} synthesises the
literature into a unified gap analysis and identifies open research directions.
Section~\ref{sec:conclusion} concludes.
